% Trang bìa, tóm tắt, mục lục. Không tính vào 15–20 trang.
\pagenumbering{roman}
\begin{titlepage}
\newgeometry{margin=2cm}
\begin{tikzpicture}[remember picture,overlay]
\draw[line width=1.2pt] ([xshift=1.5cm,yshift=-1.5cm]current page.north west) rectangle ([xshift=-1.5cm,yshift=1.5cm]current page.south east);
\draw[line width=0.5pt] ([xshift=1.65cm,yshift=-1.65cm]current page.north west) rectangle ([xshift=-1.65cm,yshift=1.65cm]current page.south east);
\end{tikzpicture}
\centering
{\large ĐẠI HỌC BÁCH KHOA HÀ NỘI}\\[4pt]
{\large\bfseries TRƯỜNG CÔNG NGHỆ THÔNG TIN VÀ TRUYỀN THÔNG}\\[1.2cm]
\includegraphics[height=4cm]{images/logo-hust.png}\\[1.2cm]
{\LARGE\bfseries Báo Cáo Chuyên Đề Nghiên Cứu}\\[1cm]
\begin{minipage}{0.88\textwidth}\centering\large
\textbf{\textit{\underline{Đề tài}:}} \textbf{Phát triển theo đặc tả so với vibe coding: nghiên cứu tình huống và thực nghiệm đối chứng trên một dự án phần mềm AI cho ngành dược Việt Nam}\\[10pt]
\textbf{\textit{\underline{Học phần}:}} \textbf{Chuyên đề nghiên cứu 1 (IT5421)}
\end{minipage}\\[1cm]
{\large Giảng viên: \textbf{TS. Nguyễn Thanh Hùng}}\\[6pt]
{\large Nhóm học viên thực hiện:}\\[6pt]
\begin{tabular}{|c|l|c|}
\hline
\textbf{STT} & \textbf{Họ và tên} & \textbf{MSHV}\\ \hline
1 & Doãn Văn An & 20252778M\\ \hline
2 & Lê Xuân Lợi & 20252781M\\ \hline
3 & Trần Ngọc Bảo & 20251227M\\ \hline
\end{tabular}
\vfill
{\large\bfseries Hà Nội, năm 2026}
\restoregeometry
\end{titlepage}

\section*{Tóm tắt}\zlabel{chk:abstract}
Phát triển theo đặc tả, tức để trợ lý lập trình AI viết thiết kế, kế hoạch và test trước khi viết mã, thường
được giới thiệu như cách bù cho điểm yếu của \emph{vibe coding}, nơi người dùng giao việc bằng vài câu ngắn.
Bằng chứng định lượng cho lời hứa này còn ít. Nghiên cứu này so sánh hai cách làm trên ViPharmind, một dự án
phần mềm AI cho ngành dược Việt Nam được xây dựng với trợ lý lập trình. Phần đầu phân tích lịch sử
git của dự án (\NumCommits{} commit). Phần thực nghiệm dựng \NumTasksMain{} nhiệm vụ từ commit thật ở ba tầng kích
thước, chấm bằng test mà agent không thấy (hidden test), và giao mỗi nhiệm vụ cho ba cohort (Claude Code với
Sonnet 5.5 và Opus 5.5, Codex) theo hai cách làm, mỗi cách ba lần, tổng \NumResRuns{} lượt. Giả thuyết, đề
bài và mã phân tích được chốt bằng git tag trước khi mở kết quả.

Khi đề bài đã nêu đủ tên và chữ ký, làm theo đặc tả không làm tăng tỉ lệ đạt. Phép thử xác nhận (Sonnet,
nhiệm vụ nhỏ) cho hiệu số \NumResHOneEstPp{} điểm phần trăm, khoảng tin cậy 95\%
[\NumResHOneLowPp{}; \NumResHOneHighPp{}], trong khi chi phí tăng gấp \NumResTokenRatioMin{} đến
\NumResTokenRatioMax{} lần token. \NumLabOutside{} trên \NumLabRuns{} lượt hỏng của Claude có nguyên nhân là
chi tiết đề không chốt hoặc nối sai với mã có sẵn, với tỉ lệ như nhau ở hai cách làm. Một thí nghiệm thăm dò
với đề chỉ một hai câu cho thấy bỏ chi tiết làm điểm hành vi giảm \NumETwoXScoreShiftSpec{} và
\NumETwoXScoreShiftVibe{} điểm phần trăm, còn quy trình đặc tả chỉ bù lại \NumETwoPooledSmallPp{} điểm ở
nhiệm vụ nhỏ. Kết quả gợi ý rằng độ rõ của đề quan trọng hơn quy trình: quy trình tự động giúp agent tự
kiểm, nhưng không tạo ra thông tin mà chỉ người giao việc có.

\textbf{Từ khóa:} coding agent, phát triển theo đặc tả, vibe coding, thực nghiệm đối chứng, preregistration.

\clearpage
\tableofcontents
